\documentclass[a4paper,12pt]{article}
\usepackage[T2A]{fontenc}
\usepackage[utf8]{inputenc}
\usepackage[russian]{babel}
\usepackage{amsmath,amssymb,mathtools}
\usepackage{helvet}
\renewcommand{\familydefault}{\sfdefault}
\usepackage{icomma}
\usepackage[a4paper,left=24mm,right=24mm,top=23mm,bottom=24mm]{geometry}
\usepackage{microtype}
\usepackage{tikz}
\usetikzlibrary{arrows.meta,positioning,shapes.geometric,calc}
\usepackage{tabularx,booktabs}
\usepackage{enumitem}
\usepackage{fancyhdr}
\usepackage{setspace}
\usepackage{xcolor}
\definecolor{accent}{HTML}{008080}
\definecolor{darkaccent}{HTML}{005B5B}
\definecolor{textgray}{HTML}{333333}
\definecolor{softgray}{HTML}{666666}
\color{textgray}
\onehalfspacing
\setlength{\parindent}{0pt}
\setlength{\parskip}{5pt}
\setcounter{secnumdepth}{1}
\usepackage[colorlinks=true,linkcolor=darkaccent]{hyperref}
\hypersetup{pdftitle={Чек-лист. Производные и исследование функции},pdfauthor={Лёвин Артём Александрович}}
\pagestyle{fancy}
\fancyhf{}
\fancyhead[L]{\small\color{softgray} Производные и исследование функции}
\fancyhead[R]{\small\color{softgray} Матвей Горбачев}
\fancyfoot[C]{\small\thepage}
\renewcommand{\headrulewidth}{0.35pt}
\renewcommand{\headrule}{\hbox to\headwidth{\color{accent}\leaders\hrule height \headrulewidth\hfill}}
\setlength{\headheight}{25pt}
\newcommand{\idea}[1]{\textcolor{darkaccent}{\textbf{#1}}}
\newcommand{\ex}[1]{\par\medskip{\color{darkaccent}\textbf{Пример. #1}}\par\smallskip}
\newcommand{\alert}[1]{\textcolor{darkaccent}{\textbf{Важно.}} #1}
\newcommand{\solution}[1]{{\small\color{softgray}#1}}
% В основных разделах используется стандартная нумерация с тонкими разделителями.
\makeatletter
\def\@seccntformat#1{\textcolor{darkaccent}{\csname the#1\endcsname}.\quad}
\makeatother
\begin{document}
\thispagestyle{empty}
\begin{center}
\vspace*{16mm}
{\large\color{darkaccent} МАТЕМАТИКА}\par
\vspace{22mm}
{\LARGE Чек-лист}\par
\vspace{15mm}
{\Large\color{accent} Производные и\\[4pt] исследование функции}\par
\vspace{20mm}
{\normalsize Подготовка к ЕГЭ по профильной математике}\par
\vfill
{\large Матвей Горбачев}\par
\vspace{13mm}
{\normalsize 5 ноября 2026 г.}\par
\vfill
{\small Лёвин Артём Александрович\\ эксклюзивно для Матвея Горбачева}\par
\vspace{15mm}
\begin{tikzpicture}[x=1cm,y=1cm]
\draw[accent,line width=0.75pt] (0,0) .. controls (1.0,0.42) and (1.5,-0.32) .. (3,0)
 .. controls (4.2,0.42) and (4.9,-0.35) .. (6,0)
 .. controls (7.1,0.37) and (7.7,-0.30) .. (9,0);
\end{tikzpicture}
\vspace{6mm}
\end{center}
\clearpage
\section{Производная: базовые правила}
\textbf{Производная} $f'(x)$ позволяет определять, как изменяется функция: возрастает она или убывает. На занятии мы начали с самых простых производных и перешли к исследованию функций без графика.

\idea{Таблица основных формул}\par
\begin{center}
\renewcommand{\arraystretch}{1.42}
\begin{tabular}{@{}ll@{}}
\toprule
Функция $f(x)$ & Производная $f'(x)$\\
\midrule
$C$ (константа) & $0$\\
$x^n$ & $nx^{n-1}$\\
$\sqrt{x}$ & $\dfrac{1}{2\sqrt{x}}$, $x>0$\\
$\sin x$ & $\cos x$\\
$\cos x$ & $-\sin x$\\
$\tg x$ & $\dfrac{1}{\cos^2 x}$, $\cos x\ne0$\\
$\ctg x$ & $-\dfrac{1}{\sin^2 x}$, $\sin x\ne0$\\
$a^x$, $a>0$, $a\ne1$ & $a^x\ln a$\\
$e^x$ & $e^x$\\
\bottomrule
\end{tabular}
\end{center}

Если $C$ --- постоянное число, то
\[
(Cf(x))'=Cf'(x),\qquad (f(x)\pm g(x))'=f'(x)\pm g'(x).
\]
\ex{Степень и постоянный множитель}
\[
\bigl(x\sqrt{x}\bigr)'=\bigl(x^{3/2}\bigr)'=\frac32\sqrt{x},\quad x>0;
\qquad (7x^3+5)'=21x^2.
\]
\alert{Постоянный множитель сохраняется.} Производная суммы считается по слагаемым. Например, $(7x^2)'=14x$, но $(7+x^2)'=2x$.

\section{Произведение и частное}
Для двух функций $u(x)$ и $v(x)$:
\[
\boxed{(uv)'=u'v+uv'},\qquad
\boxed{\left(\frac uv\right)'=\frac{u'v-uv'}{v^2}},\quad v\ne0.
\]
\idea{Рабочая запись:} сначала выпишите $u$, $v$, $u'$, $v'$, затем подставьте их в формулу. Это снижает риск перепутать множители и знаки.

\ex{Произведение}
\[
 y=(x^2+4)(x^3-5x).
\]
\[
\begin{aligned}
u&=x^2+4, & u'&=2x,\\
v&=x^3-5x, & v'&=3x^2-5;
\end{aligned}
\]
\[
 y'=2x(x^3-5x)+(x^2+4)(3x^2-5).
\]
\ex{Частное с теми же функциями}
\[
\left(\frac{x^2+4}{x^3-5x}\right)'=
\frac{2x(x^3-5x)-(x^2+4)(3x^2-5)}{(x^3-5x)^2}.
\]
Область определения исходной дроби: $x\ne0$, $x\ne\sqrt5$, $x\ne-\sqrt5$.

\section{Производная сложной функции}
Если $y=F(g(x))$, то
\[
\boxed{y'=F'(g(x))\cdot g'(x)}.
\]
\idea{Как распознать:} внутри известной функции вместо простого $x$ стоит выражение. Сначала найдите производную внешней функции, сохранив внутреннее выражение, затем умножьте на производную внутреннего.

\ex{Степень скобки}
\[
\begin{aligned}
y&=(5x-3)^{10},\\
y'&=10(5x-3)^9\cdot 5=50(5x-3)^9.
\end{aligned}
\]
\ex{Косинус сложного аргумента}
\[
\begin{aligned}
y&=\cos(4x^2-10),\\
y'&=-\sin(4x^2-10)\cdot 8x.
\end{aligned}
\]
\ex{Показательная функция}
\[
\begin{aligned}
y&=2^{x^2-x-6},\\
y'&=2^{x^2-x-6}\ln 2\cdot(2x-1).
\end{aligned}
\]
Для $e^{-3x+5}$ производная равна $-3e^{-3x+5}$.

\alert{Три частые ошибки:} забыть производную внутреннего выражения; ошибиться в знаке производной косинуса; заменить $(5x-3)^9$ на $5^9$.

\section{Что показывает знак производной}
На каждом интервале, где функция непрерывна и дифференцируема:
\[
\begin{array}{lll}
f'(x)>0 &\Longrightarrow& f(x)\ \text{возрастает};\\
f'(x)<0 &\Longrightarrow& f(x)\ \text{убывает}.
\end{array}
\]
Если при переходе через точку знак $f'$ меняется с минуса на плюс, получаем \textbf{локальный минимум}; если с плюса на минус --- \textbf{локальный максимум}. Одного равенства $f'(x)=0$ недостаточно, чтобы утверждать наличие экстремума.

\ex{Исследование квадратичной функции}
\[
 f(x)=x^2-x-6,\quad f'(x)=2x-1.
\]
Решаем $2x-1=0$: $x=\frac12$. Слева $f'<0$, справа $f'>0$. Значит, функция убывает при $x<\frac12$ и возрастает при $x>\frac12$.
\[
 f\!\left(\frac12\right)=\frac14-\frac12-6=-\frac{25}{4}.
\]
Итак, точка минимума --- $x=\frac12$, минимальное значение --- $-\frac{25}{4}$.

\ex{Два стационарных значения}
\[
 f(x)=\frac{x^3}{3}-\frac{x^2}{2}-6x+7,
\quad f'(x)=x^2-x-6=(x+2)(x-3).
\]
Знаки $f'$ по интервалам $(-\infty,-2)$, $(-2,3)$, $(3,+\infty)$ равны соответственно $+$, $-$, $+$. Следовательно, $x=-2$ --- точка локального максимума, $x=3$ --- локального минимума.

\alert{Проверяйте знаки на каждом интервале подстановкой пробного значения.} При переходе через нуль производной знак может и не измениться: например, для $f'(x)=3x^2$ он положителен с обеих сторон от нуля.

\clearpage
\section*{Алгоритм: исследование функции по производной}
\vspace{4mm}
\begin{center}
\tikzset{
  flowstep/.style={draw=accent,rounded corners=2pt,line width=.65pt,minimum width=11.7cm,minimum height=1.05cm,text width=11.1cm,align=center,inner sep=5pt,font=\normalsize},
  decision/.style={draw=accent,diamond,aspect=3.0,line width=.65pt,align=center,text width=3.5cm,inner sep=1pt,font=\normalsize},
  branch/.style={draw=accent,rounded corners=2pt,minimum width=5.1cm,minimum height=1.05cm,align=center,font=\normalsize},
  arr/.style={-{Latex[length=2.3mm]},line width=.7pt,color=darkaccent}
}
\begin{tikzpicture}[node distance=7mm]
\node[flowstep] (a) {1. Найдите область определения $D(f)$ и исключённые точки.};
\node[flowstep,below=of a] (b) {2. Вычислите $f'(x)$ и укажите, где она существует.};
\node[flowstep,below=of b] (c) {3. Найдите корни $f'(x)=0$ и точки из $D(f)$, где $f'$ не существует.};
\node[flowstep,below=of c] (d) {4. Отметьте все найденные точки и разрывы области определения на оси.};
\node[flowstep,below=of d] (e) {5. На каждом полученном интервале подставьте пробное $x$ \emph{в производную}.};
\node[decision,below=8mm of e] (q) {Каков знак $f'(x)$?};
\node[branch,below left=10mm and 5mm of q] (l) {$f'<0$\\функция убывает};
\node[branch,below right=10mm and 5mm of q] (r) {$f'>0$\\функция возрастает};
\node[flowstep,below=24mm of q] (z) {6. По смене знаков найдите экстремумы: $+\to-$ --- максимум, $-\to+$ --- минимум. Затем вычислите значения функции.};
\draw[arr] (a)--(b); \draw[arr](b)--(c); \draw[arr](c)--(d); \draw[arr](d)--(e); \draw[arr](e)--(q);
\draw[arr] (q.west) -| node[pos=.22,above,font=\small]{минус} (l.north);
\draw[arr] (q.east) -| node[pos=.22,above,font=\small]{плюс} (r.north);
\draw[arr] (l.south) -- ([xshift=-3cm]z.north);
\draw[arr] (r.south) -- ([xshift=3cm]z.north);
\end{tikzpicture}
\end{center}
\clearpage
\section{Дробно-рациональная функция: важные ограничения}
При постоянном \textbf{положительном} числителе и ненулевом знаменателе $g(x)$ дробь
\[
 f(x)=\frac{C}{g(x)},\quad C>0,
\qquad f'(x)=-\frac{C\,g'(x)}{g(x)^2}
\]
на каждом интервале своей области определения меняется в направлении, противоположном знаменателю. Этот приём позволяет исследовать функцию без громоздких вычислений.

\ex{Функция с разрывами}
\[
 f(x)=\frac{2}{x^2-4},\qquad D(f)=\mathbb R\setminus\{-2,2\}.
\]
\[
 f'(x)=-\frac{4x}{(x^2-4)^2}.
\]
Знаменатель производной положителен везде, где функция определена. Поэтому $f'(x)>0$ при $x<0$ и $f'(x)<0$ при $x>0$ с учётом исключённых точек $\pm2$.

\begin{center}
\small
\renewcommand{\arraystretch}{1.35}
\begin{tabular}{@{}lll@{}}
\toprule
Интервал & Знак $f'(x)$ & Поведение функции\\
\midrule
$(-\infty,-2)$, $(-2,0)$ & $+$ & возрастает\\
$(0,2)$, $(2,+\infty)$ & $-$ & убывает\\
\bottomrule
\end{tabular}
\end{center}

В точке $x=0$ производная меняет знак с плюса на минус, следовательно, здесь \textbf{локальный максимум}: $f(0)=-\frac12$. Глобального максимума у функции нет: при $x\to2+0$ значение $f(x)$ неограниченно растёт.

\alert{Нельзя} объединять интервалы через точку, исключённую из ОДЗ. Если в числителе переменная или отрицательная константа, правило о противоположном поведении знаменателя в указанном виде неприменимо.

\section{Короткая памятка перед решением}
\begin{enumerate}[leftmargin=*,itemsep=1pt,topsep=0pt]
\item Сначала найдите ОДЗ и отделите переменные от постоянных.
\item Выберите правило: сумма, произведение, частное или сложная функция.
\item При необходимости выпишите $u$, $v$, $u'$, $v'$ и производную внутреннего выражения.
\item Исследуя функцию, проверьте знак $f'$ на каждом интервале области определения.
\end{enumerate}

\clearpage
\section*{Тренировочный блок}\label{training}
\textcolor{darkaccent}{Путь А --- 10 заданий.} Выполняйте их последовательно. В заданиях 9--10 требуется полный анализ знаков производной с указанием интервалов и экстремумов.
\vspace{4mm}
\begin{enumerate}[leftmargin=*,itemsep=9pt]
\item Найдите производную $y=7x^4-3x^2+9$.
\item Найдите производную $y=x\sqrt{x}+\dfrac4x$, $x>0$.
\item Найдите производную $y=(x^2-3)(x^3+2x)$.
\item Найдите производную $y=\dfrac{2x+1}{x^2+1}$.
\item Найдите производную $y=\sin x\cos x$.
\item Найдите производную $y=(3x-2)^7$.
\item Найдите производную $y=\cos(2x^2-5x)$.
\item Найдите производную $y=5^{x^2-2x}$.
\item Исследуйте $f(x)=x^3-3x$: найдите интервалы возрастания и убывания, точки и значения локальных экстремумов.
\item Исследуйте $f(x)=\dfrac{2}{x^2-9}$: определите ОДЗ, интервалы возрастания и убывания, точки и значения локальных экстремумов.
\end{enumerate}
\vfill
\textcolor{softgray}{\small Для задач с дробями обязательно указывайте исключённые значения переменной. Интервалы записывайте с соблюдением ОДЗ.}
\clearpage
\section*{Ответы к тренировочному блоку}
\renewcommand{\arraystretch}{1.6}
\begin{tabularx}{\linewidth}{@{}p{1cm}X@{}}
\toprule
№ & Ответ\\
\midrule
1 & $y'=28x^3-6x$.\\
2 & $y'=\dfrac32\sqrt{x}-\dfrac4{x^2}$, $x>0$.\\
3 & $y'=5x^4-3x^2-6$.\\
4 & $y'=\dfrac{-2x^2-2x+2}{(x^2+1)^2}$.\\
5 & $y'=\cos^2x-\sin^2x=\cos2x$.\\
6 & $y'=21(3x-2)^6$.\\
7 & $y'=-(4x-5)\sin(2x^2-5x)$.\\
8 & $y'=(2x-2)5^{x^2-2x}\ln5$.\\
9 & $f'=3(x-1)(x+1)$. Возрастает: $(-\infty,-1)$ и $(1,+\infty)$; убывает: $(-1,1)$. Локальный максимум: $f(-1)=2$; локальный минимум: $f(1)=-2$.\\
10 & $D(f)=\mathbb R\setminus\{-3,3\}$; $f'=-\dfrac{4x}{(x^2-9)^2}$. Возрастает: $(-\infty,-3)$ и $(-3,0)$; убывает: $(0,3)$ и $(3,+\infty)$. Локальный максимум: $f(0)=-\dfrac29$; локального минимума нет.\\
\bottomrule
\end{tabularx}
\vfill
\begin{center}
{\small\color{softgray} Лёвин Артём Александрович\\ эксклюзивно для Матвея Горбачева}
\end{center}
\end{document}
